\section*{अध्याय \ref{ch:g12:ph-conductivity-titrations} --- pH और चालकता द्वारा अनुमापन}

\begin{solution}{exo:g12:ph-conductivity-titrations:1}
$V_E = \qty{20.0}{mL}$; $c = 0.100 \times 20.0 / 20.0 = \qty{0.100}{mol/L}$।
\end{solution}

\begin{solution}{exo:g12:ph-conductivity-titrations:2}
\omterm{def:g12:ph-conductivity-titrations:ph-metric}{अर्ध-तुल्यता} पर \omterm{def:g9:acids-bases-ph:ph}{pH} $= pK_a$: $pK_a = 3.75$, मेथेनोइक अम्ल
(3.74) के निकट।
\end{solution}

\begin{solution}{exo:g12:ph-conductivity-titrations:3}
फ़ीनॉल्फ़्थेलीन: उसका परास 8.0--10.0 उछाल में आता है, जो क्षारकीय है।
मेथिल रेड (4.2--6.3) \omterm{def:g12:ph-conductivity-titrations:ph-metric}{अर्ध-तुल्यता} के निकट ही बदल जाता, बहुत जल्दी।
\end{solution}

\begin{solution}{exo:g12:ph-conductivity-titrations:4}
प्रोब की अनुक्रिया समय और ताप के साथ खिसकती है: उसे ज्ञात \omterm{def:g9:acids-bases-ph:ph}{pH} वाले
दो \omterm{def:g12:buffers-predominance:buffer}{बफ़र विलयनों} में सही पढ़ने के लिए समायोजित किया जाता है।
\end{solution}

\begin{solution}{exo:g12:ph-conductivity-titrations:5}
\ce{H3O+} (349.8) > \ce{OH-} (198.3) > \ce{Cl-} (76.2) > \ce{Na+} (50.3)।
\end{solution}

\begin{solution}{exo:g12:ph-conductivity-titrations:6}
$c = 0.050 \times 14.6 / 20.0 = \qty{0.0365}{mol/L}$।
\end{solution}

\begin{solution}{exo:g12:ph-conductivity-titrations:7}
प्रवणताएँ: 2.0, 4.5, 17.5 और \qty{2.5}{mL^{-1}}। सबसे बड़ी
10.0 और \qty{10.2}{mL} के बीच है: $V_E \approx \qty{10.1}{mL}$।
\end{solution}

\begin{solution}{exo:g12:ph-conductivity-titrations:8}
\omterm{def:g11:titration:equivalence}{तुल्यता} से पहले मिलाया गया हर \ce{OH-} एक \ce{H3O+} (अत्यधिक
चालक) हटाता है और एक \ce{Na+} (कहीं कम चालक) लाता है:
\omterm{def:g12:ph-conductivity-titrations:conductivity}{चालकता} घटती है। उसके बाद \ce{Na+} और \ce{OH-} बिना अभिक्रिया किए
जमा होते जाते हैं: वह बढ़ती है।
\end{solution}

\begin{solution}{exo:g12:ph-conductivity-titrations:9}
आरंभ में कम \omterm{def:g9:ions:ion}{आयन} (\omterm{def:g12:ka-and-pka:strong-acid}{दुर्बल अम्ल} पानी से मुश्किल से अभिक्रिया करता है): कम
\omterm{def:g12:ph-conductivity-titrations:conductivity}{चालकता}। \omterm{def:g11:titration:equivalence}{तुल्यता} से पहले अनावेशित \omterm{def:g7:atoms-and-molecules:molecule}{अणुओं} की जगह एथेनोएट और सोडियम \omterm{def:g9:ions:ion}{आयन}
बनते हैं: वह बढ़ती है, धीरे-धीरे। \omterm{def:g11:titration:equivalence}{तुल्यता} के बाद
\ce{Na+} और \ce{OH-} जमा होते हैं: वह तेज़ी से बढ़ती है। दो बढ़ती रेखाएँ,
दूसरी अधिक खड़ी।
\end{solution}

\begin{solution}{exo:g12:ph-conductivity-titrations:10}
लगभग 2.9, 4.76 और 8.7। एथेनोइक अम्ल दुर्बल है: उसका केवल कुछ प्रतिशत
\omterm{def:g12:ka-and-pka:oxonium}{ऑक्सोनियम आयन} देता है, इसलिए \omterm{def:g9:acids-bases-ph:ph}{pH} मिलती-जुलती सांद्रता वाले हाइड्रोक्लोरिक
अम्ल के 1.0 से ऊँचा शुरू होता है।
\end{solution}

\begin{solution}{exo:g12:ph-conductivity-titrations:11}
ताकि मिलाया गया \omterm{def:g11:titration:titration}{अनुमापक} का आयतन कुल आयतन को मुश्किल से बदले: तब
सांद्रताएँ, और इसलिए \omterm{def:g12:ph-conductivity-titrations:conductivity}{चालकता}, सीधी रेखाओं के साथ बदलती हैं।
\end{solution}

\begin{solution}{exo:g12:ph-conductivity-titrations:12}
हाइड्रॉक्साइड \omterm{def:g9:ions:ion}{आयन} पहले \omterm{def:g12:ka-and-pka:strong-acid}{प्रबल अम्ल} (\ce{H3O+}) से अभिक्रिया करते हैं, फिर
\omterm{def:g12:ka-and-pka:strong-acid}{दुर्बल अम्ल} से। पहला उछाल हाइड्रोक्लोरिक अम्ल के अंत को चिह्नित करता है
(उसका \omterm{def:g11:titration:equivalence}{तुल्यता आयतन} उसे मापता है), दूसरा एथेनोइक
अम्ल के अंत को (दोनों उछालों के बीच का आयतन उसे मापता है)।
\end{solution}

\begin{solution}{exo:g12:ph-conductivity-titrations:13}
आरंभ: $(349.8 + 76.2) \times 0.0100 = \qty{4.26}{mS/cm}$।
\omterm{def:g11:titration:equivalence}{तुल्यता} पर: $[\ce{Na+}] = [\ce{Cl-}] = 1.00/110 = \qty{0.00909}{mol/L}$,
इसलिए $(50.3 + 76.2) \times 0.00909 = \qty{1.15}{mS/cm}$। दोनों
चित्र के अनुसार।
\end{solution}

\begin{solution}{exo:g12:ph-conductivity-titrations:14}
$V_E$ पर कोई नहीं: पाठ्यांकों का विचलन चाहे जो हो, उछाल उसी आयतन पर
होता है। \omterm{def:g12:ph-conductivity-titrations:ph-metric}{अर्ध-तुल्यता} पर पढ़ा गया $pK_a$ 0.2 अधिक होता है।
\end{solution}

\begin{solution}{exo:g12:ph-conductivity-titrations:15}
इस \omterm{def:g9:ions:precipitate}{अवक्षेपण} के दौरान \omterm{def:g9:acids-bases-ph:ph}{pH} नहीं बदलता, पर \omterm{def:g9:ions:ion}{आयन} बदलते हैं:
\omterm{def:g11:titration:equivalence}{तुल्यता} से पहले मिलाया गया हर \ce{Ag+} एक \ce{Cl-} हटाता है और समान \omterm{def:g12:ph-conductivity-titrations:conductivity}{चालकता} का
एक \ce{NO3-} लाता है, इसलिए \omterm{def:g12:ph-conductivity-titrations:conductivity}{चालकता} लगभग
स्थिर रहती है (सोडियम \omterm{def:g9:ions:ion}{आयन} शुरू से ही मौजूद थे); \omterm{def:g11:titration:equivalence}{तुल्यता} के बाद \ce{Ag+} और \ce{NO3-}
जमा होते हैं और वह बढ़ती है। एक समतल रेखा, फिर एक बढ़ती रेखा।
\end{solution}

\begin{solution}{pb:g12:ph-conductivity-titrations:1}
\textbf{1.} सिरके के प्रति \qty{100}{g} एथेनोइक अम्ल के \qty{6}{g}।

\textbf{2.} सिरके के \qty{101}{g} में $6.0 \times 101/100 =
\qty{6.06}{g}$ होते हैं; $M = \qty{60.0}{g/mol}$: \qty{0.1000}{L} में
\qty{0.101}{mol}, अर्थात \qty{1.01}{mol/L}।

\textbf{3.} बिना तनु किए, \qty{10.0}{mL} के लिए \qty{0.100}{mol/L} सोडियम हाइड्रॉक्साइड के
लगभग \qty{101}{mL} चाहिए होते: ब्यूरेट की क्षमता से अधिक।

\textbf{4.} \qty{10.0}{mL} का आयतनमितीय पिपेट और \qty{100.0}{mL} का
आयतनमितीय फ़्लास्क।

\textbf{5.} \ce{CH3COOH + OH- -> CH3COO- + H2O}।

\textbf{6.} $V_E = \qty{10.1}{mL}$।

\textbf{7.} $c = 0.100 \times 10.1 / 10.0 = \qty{0.101}{mol/L}$।

\textbf{8.} $10 \times 0.101 = \qty{1.01}{mol/L}$।

\textbf{9.} \qty{5.05}{mL} पर \omterm{def:g9:acids-bases-ph:ph}{pH} 4.76: एथेनोइक अम्ल का $pK_a$।

\textbf{10.} \omterm{def:g11:titration:equivalence}{तुल्यता} पर विलयन में सोडियम एथेनोएट होता है, और
एथेनोएट एक \omterm{def:g12:ka-and-pka:strong-acid}{दुर्बल क्षारक} है।

\textbf{11.} फ़ीनॉल्फ़्थेलीन, जिसके परास 8.0--10.0 में \omterm{def:g11:titration:equivalence}{तुल्यता} का \omterm{def:g9:acids-bases-ph:ph}{pH}
(लगभग 8.7) आता है।

\textbf{12.} सोडियम \omterm{def:g9:ions:ion}{आयन} और एथेनोएट \omterm{def:g9:ions:ion}{आयन}।

\textbf{13.} मिलाया गया हर \ce{OH-} लगभग अनआयनित एक \omterm{def:g7:atoms-and-molecules:molecule}{अणु} को
\ce{CH3COO-} में बदलता है, साथ में एक \ce{Na+}: साधारण \omterm{def:g12:ph-conductivity-titrations:conductivity}{चालकता} वाले \omterm{def:g9:ions:ion}{आयन}
धीरे-धीरे जुड़ते हैं।

\textbf{14.} \omterm{def:g11:titration:equivalence}{तुल्यता} के बाद अत्यधिक चालक हाइड्रॉक्साइड \omterm{def:g9:ions:ion}{आयन}
सोडियम \omterm{def:g9:ions:ion}{आयनों} के साथ जमा होते हैं।

\textbf{15.} \omterm{def:g12:ka-and-pka:strong-acid}{दुर्बल अम्ल} पानी को बहुत कम \omterm{def:g9:ions:ion}{आयन} देता है।

\textbf{16.} $V_E = \qty{10.1}{mL}$ पर, जैसा \omterm{def:g12:ph-conductivity-titrations:ph-metric}{pH-मितीय अनुमापन} में।

\textbf{17.} $1.01 \times 60.0 = \qty{60.6}{g/L}$।

\textbf{18.} \qty{100.0}{mL} में \qty{6.06}{g}।

\textbf{19.} प्रति \qty{100}{g} $6.06 \times 100 / 101 = \qty{6.0}{g}$।

\textbf{20.} सिरके की अम्लता \qty{6.0}{\percent} है: लेबल
सही है।
\end{solution}
